\section{Mark Reinhold}
Chief Architect do Java Platform Group na Oracle, onde trabalha no Java~SE e
no OpenJDK (fonte: \citeobra{bio:reinhold-inside}).

\begin{itemize}
\obra{doc:sotms}{O sistema de módulos explicado em uma leitura:
  \texttt{module-info}, readability, accessibility e módulos automáticos.}
\obra{video:reinhold-modules-one-lesson}{Os módulos ao vivo no terminal,
  em uma aula.}
\obra{doc:jep261}{A especificação de implementação do sistema de módulos
  (Owner: Reinhold).}
\obra{doc:jep403}{Por que o \texttt{--illegal-access} acabou e o que isso
  quebra em bibliotecas antigas que usam reflexão.}
\obra{video:reinhold-unsafe}{A história do \texttt{sun.misc.Unsafe} e por
  que ele vai sumindo.}
\obra{artigo:forward-faster}{A proposta de 2017 que criou a cadência de
  uma versão a cada seis meses.}
\obra{artigo:forward-even-faster}{A proposta de 2021 de uma LTS a cada dois
  anos.}
\obra{doc:jep3}{O processo de lançamento do JDK, por quem o desenhou.}
\obra{doc:leyden-beginnings}{Os objetivos do Project Leyden (startup,
  warmup, footprint): bom contraponto ao build time do Quarkus.}
\end{itemize}

\section{Paul Sandoz}
Architect no Java Platform Group da Oracle (fonte:
\citeobra{bio:sandoz-inside}). A página tem pouca bio; os JEPs e projetos
abaixo, que ele lidera, mostram o papel dele.

\begin{itemize}
\obra{doc:jep193}{VarHandles: acesso atômico e com ordenação de memória,
  a alternativa segura ao \texttt{Unsafe} (Owner: Sandoz; autor: Doug Lea).}
\obra{doc:jep338}{A Vector API: SIMD em Java puro, compilado pelo C2 para
  instruções vetoriais. Passou por várias rodadas de incubação, à espera do Valhalla.}
\obra{video:sandoz-vector-api}{Demonstração da Vector API e do código de
  máquina gerado, com Sandhya Viswanathan, da Intel.}
\obra{doc:babylon-code-models}{Code reflection: representar código Java
  simbolicamente para transformá-lo (GPU, SQL, diferenciação automática).}
\obra{video:sandoz-babylon}{O Project Babylon explicado pelo líder do
  projeto.}
\obra{video:sandoz-java-ai}{FFM, Vector API, Valhalla e Babylon aplicados a
  IA (JavaOne 2025).}
\obra*{doc:leyden-condensers}{Os \enquote{condensadores} do Leyden:
  transformar a aplicação antes de rodar, ideia parente do build time do
  Quarkus.}
\end{itemize}
